\input{../tex-inputs/latex-leseansicht-vorspann}

\section[Arthur Schnitzler an Gustav Schwarzkopf, {[}zwischen 1897 und Mitte 1898?{]}]{L04189 Arthur Schnitzler an Gustav Schwarzkopf, {[}zwischen 1897 und Mitte 1898?{]}}
\nopagebreak\mylabel{L04189v}
\rehead{ }\normalsize\beginnumbering\briefempfaengerindex{Schwarzkopf, Gustav@\textsc{Schwarzkopf, Gustav}!zzzSchnitzler, Arthur@\emph{von Arthur Schnitzler}!1898-06-301@{{[}zwischen 1897 und Mitte 1898?{]}}|(be}
\toendnotes[C]{\smallbreak\pagebreak[2]}
\correspDesc{Versand  durch Arthur Schnitzler im Zeitraum [zwischen 1897 und Mitte 1898?] in Wien
\newline{}Erhalt  durch Gustav Schwarzkopf im Zeitraum [zwischen 1897 und Mitte 1898?] in Wien}\toendnotes[C]{\smallbreak}
\Standort{CUL, Schnitzler, B 96.}
\physDesc{Brief, 1 Blatt, 2 Seiten, 104 Zeichen
\newline{}Handschrift: Bleistift, deutsche Kurrent}\toendnotes[C]{\smallbreak}
\pstart{}{\pb}\label{K_L04189-1v}\edtext{Lieber Guſtav}{\lemma{\textnormal{\emph{Lieber Gustav}}}\Cendnote{\textnormal{Das Blatt ist undatiert. Durch die Verwendung der Anrede »Lieber Guſtav« dürfte
                     es frühestens 1897 übermittelt sein. Das häufige Frequentieren des Café Arkaden\oindex{Wien@\textbf{Wien}!I., Innere Stadt@\textbf{I., Innere Stadt}!Café Arkaden@\textbf{Café Arkaden}|pwk} endete spätestens im Frühjahr 1898, womit eine hintere 
                  zeitliche Grenze vorzuliegen scheint.}}}\label{K_L04189-1},\pend\vspace{0.5em}
\pstart
           bitte mich nicht abzuholen. Vielleicht ſehen {\pb}wir uns heut Abd im \textsc{Arkaden\oindex{Wien@\textbf{Wien}!I., Innere Stadt@\textbf{I., Innere Stadt}!Café Arkaden@\textbf{Café Arkaden}|pw}}?\pend
           
\pstart
           Herzlich{\\[\baselineskip]} Ihr \spacefill\mbox{Arth Sch}\pend
           \leftskip=0em{}\selectlanguage{ngerman}\endnumbering\briefempfaengerindex{Schwarzkopf, Gustav@\textsc{Schwarzkopf, Gustav}!zzzSchnitzler, Arthur@\emph{von Arthur Schnitzler}!1897-01-011@{{[}zwischen 1897 und Mitte 1898?{]}}|)be}\mylabel{L04189h}
\begin{anhang}
\end{anhang}\newcommand{\dateiname}{L04189}\newcommand{\titel}{Arthur Schnitzler an Gustav Schwarzkopf, [zwischen 1897 und Mitte 1898?]}\newcommand{\editorInnen}{Herausgegeben von Jahnke, SelmaMüller, Martin Anton}\input{../tex-inputs/latex-leseansicht-abspann}
